\section{The finite laws and their continuum observables}\label{sec:regularization}

We use the all-flux chamber model of \cite[Section~2]{ChamberLattice}.
The following definitions fix its normalization and its relation to the
physical scales in \eqref{eq:cutoff-scales}. All cochain norms and
pairings sum over positively oriented cells. Reversing the orientation
changes the sign of a cochain.

\subsection{The chamber matrix and the flux sum}
Let $O_{\rm cub}$ be the signed permutation matrices in three dimensions.
A chamber label is $(c,g)$, where $c\in\Z^3$ and $g\in O_{\rm cub}$.
There are $48$ labels at each photon site. Let $s_0$ reverse coordinate
three, let $s_1,s_2$ exchange coordinates $(2,3)$ and $(1,2)$, and let
$r_{e_1}$ reverse coordinate one. With the Pauli matrices $\sigma_j$, put
$\gamma_j=\operatorname{diag}(\sigma_j,-\sigma_j)$ and
\[
 \begin{gathered}
 (w_0,w_1,w_2,w_3)=(6,6\sqrt2,6\sqrt2,6),\\
 (\mathbf n_0,\mathbf n_1,\mathbf n_2,\mathbf n_3)
 =(e_3,(e_2-e_3)/\sqrt2,(e_1-e_2)/\sqrt2,-e_1),\qquad
 M_i(g)=\tfrac12w_i(\rho+\gamma\cdot g\mathbf n_i).
 \end{gathered}
\]
For a bond phase $B$, the unit-mesh massless matrix is
\begin{equation}\label{eq:reg-dirac}
 \begin{split}
 (D^{\rm ch}[B]\psi)(c,g)={}&
 \sum_{i=0}^2M_i(g)[\psi(c,g)-\psi(c,gs_i)]\\
 &+M_3(g)[\psi(c,g)-e^{iB(c,c+ge_1)}
                         \psi(c+ge_1,gr_{e_1})].
 \end{split}
\end{equation}
On $C(n)=\Z^3/n\Z^3$ the bond fields are periodic and spinors
satisfy $\psi(c+ne_j,g)=-\psi(c,g)$. The massive matrix
$D^{\rm ch}[B]+\boldsymbol\mu$ is the matrix in
\cite[Section~2]{ChamberLattice} at diagonal parameter
$\boldsymbol\mu+\frac\rho2\sum_iw_i$, where $\boldsymbol\mu$ is
the total lattice mass. The chamber inner product has weight $1/48$;
this constant weight does not change a finite-dimensional determinant.

For $q>0$, put $r_q=2\pi/q$. A photon configuration is a pair
$(A,z)$ with $A$ a real one-cochain and $z$ a closed integer
two-cochain, modulo
\[
 (A,z)\sim(A+r_q k,z-dk),\qquad k\text{ an integer one-cochain}.
\]
Its curvature is $u=dA+r_qz$. With the representative
$A(b)\in[0,r_q)$, define
\begin{equation}\label{eq:reg-functional}
 Z_{C(n)}[G]=
 \sum_{\substack{z\in C^2(C(n);\Z)\\dz=0}}
 \int_{[0,r_q)^{3n^3}}
 G(A,z)e^{-\|u\|^2/2}
 \det(D^{\rm ch}[qA]+\boldsymbol\mu)\,\dd A.
\end{equation}
The determinant is nonnegative and $0<Z_{C(n)}[1]<\infty$ for
every real $\boldsymbol\mu$, by
\cite[Propositions~2.2--2.3]{ChamberLattice}.
We write $\E_n^{\rm af}=Z_{C(n)}[\,\cdot\,]/Z_{C(n)}[1]$.
Rectangular tori carry the same periodic photon and antiperiodic
fermion conventions. No flux class is excluded from this law.

At the physical cutoff, $q=e\sqrt\eta$ and
$\boldsymbol\mu=\boldsymbol\mu_N=\bar m\eta+m_{0,N}(e)$.
The physical field $F^\sharp$ is the centered step interpolation
of $\eta^{-3/2}u$, defined precisely in \cref{sec:continuum-limits}.
For a fixed mass sequence, $P_{N,T}$ denotes its law on $H_{-4}$.
Let $\mathcal V_N$ be the set of its weak limits as $T\to\infty$.
A volume-first limit is a weak limit of
$\mu_i\in\mathcal V_{N_i}$ with $N_i\to\infty$.

\subsection{Gaussian bounds and reflection quotients}\label{sec:probability}
For real Schwartz two-forms, Gaussian domination means
$\E_\lambda e^{F(f)}\le e^{\|f\|_2^2/2}$.
The estimates of \cite[Proposition~2.1]{ContinuumLimits} give
\[
 \E_\lambda F(f)=0,\qquad
 \E_\lambda|F(f)|^p\le c_p\|f\|_2^p,\qquad
 c_p=2\,2^{p/2}\Gamma(1+p/2).
\]
For complex tests the product bound has the additional factor
$2^{n/2}$ at degree $n$. Applied to differences, these estimates
extend $F$ from Schwartz tests to $L^2$, consistently in every
$L^p(\lambda)$, $p\ge1$. Expanding one factor at a time likewise
extends each joint moment as a bounded multilinear form on complex
$L^2$. These completions will be used for the plate insertions.
The same proposition gives entire finite-dimensional generating
functions, so their moments determine the finite-dimensional laws.

Write $h_\kappa=h_k e_{ij}$ for the real Hermite two-forms,
$k\in\Z_{\ge0}^3$, $i<j$, with eigenvalues
$\lambda_\kappa=2|k|_1+4$ for $1+|y|^2-\Delta$. Set
\[
 \|F\|_{-s}^2=\sum_\kappa\lambda_\kappa^{-2s}|F(h_\kappa)|^2,
 \qquad H_{-s}=\{F\in\mathcal S':\|F\|_{-s}<\infty\}.
\]
All weak convergence of laws on $H_{-4}$ uses its norm topology.
The Hermite compactness argument of \cite[Section~4]{ContinuumLimits}
shows that a closed bounded $H_{-3}$ ball is compact in $H_{-4}$.
The tensor estimate \cite[Proposition~6.1]{ContinuumLimits} gives a
continuous map $Q_\Psi\in L^2(\lambda)$ for $\Psi\in\mathcal S_n$,
equal to $\prod_aF(f_a)$ on elementary tensors, and the bound
\eqref{eq:main-growth} for its expectation.

For open $O$, let $\Sigma_O$ be generated by $F(f)$ with real
$f\in C_c^\infty(O;\Lambda^2)$. For a Euclidean motion $g$,
write $\widehat g A(F)=A(g^{-1}\cdot F)$ and $U(a)=\widehat{\tau_a}$.
Let $r_j$ reflect in $x_j=0$ and put
$\mathcal R_j=\widehat r_j$, $\Theta_j A=\overline{\mathcal R_jA}$.
For a translation-invariant law that is invariant and reflection
positive at this plane, the positive form
\[
 b_j(A,B)=\E_\lambda[(\Theta_jA)B],\qquad
 A,B\in L^2(\Sigma_{\{x_j>0\}},\lambda),
\]
has quotient completion $\mathcal H_j$, quotient map $[A]_j$,
and vector $\Omega_j=[1]_j$. The reconstruction and spectral cone
are \cite[Propositions~2.1 and~3.1]{ErgodicComponents}:
\[
 e^{-tH_j}V_j(y)[A]_j=[U(te_j+y)A]_j,\qquad
 t\ge0,\quad y\perp e_j,\qquad H_j\ge|\mathsf P_i|\quad(i\ne j).
\]
The cone conclusion uses invariance and reflection positivity at
$x_j+x_i=0$. These are the precise reconstruction inputs used in
\cref{sec:rotations}.

\subsection{Block scales and the selected pair}\label{subsec:integration-block-data}
The block construction uses an odd $L$, an integer power $M$ of $L$,
and the primitive and activity conventions of
\cite{MaxwellBlock,FermionBlock,NormalizedContraction,NonlinearStep,CoverTuning}.
Their constants will be chosen in \cref{sec:ordered-constants}.
For the reference mass in \cref{thm:main}, write
\[
 K=N+\log_Lh,\qquad a_j=L^{j-N},\quad e_j=e\sqrt{a_j},
 \quad\mu_j=\bar m a_j,\quad v_j=-\log e_j.
\]
Let $r_j$ be the least nonnegative integer power of $L$ at least
$v_j^2$. Put
\[
 R_{\rm top}=hMr_K,\qquad R^*=R_{\rm top}/L,\qquad
 \mathcal T(e)=\{hL^t:t\in\Z_{\ge0},\ hL^t\le R_{\rm top}\}.
\]
For $R\in\mathcal T(e)$, the stop is the first level
$k_*(R)$ with $\ell_j=Mr_j$, where $\ell_j=R/a_j$.
The finite threshold $N_{\min}(R,e)$ and the existence and
properties of the stop are given in \cite[Proposition~6.1]{CoverTuning}.
A star on a running quantity means evaluation at this stop.
Write $N_{\rm ref}(e)=N_{\min}(R^*,e)$ and $n_*=R^*/\eta$.

Under the numerical conditions verified in \cref{prop:ordered-constants}, the common-pair theorem
of \cite[Theorem~A]{CoverTuning} selects
$P_N(e)=(m_{0,N}(e),\varepsilon^v_{0,N}(e))$ on the top torus,
with zero terminal mass shift and vacuum term there. The same pair
is used on the smaller reference tori, for every twist and source.
Its scalar bounds and controlled trajectories are conclusions of
that theorem. Below the top definition threshold, only
$m_{0,N}(e)$ is set to zero. No trajectory is used at those cutoffs.
The construction of this pair requires no choice of the probe.

The early zero-flux source identities in
\cite[Section~3]{ReferenceResponse} define the matter normalization
used by \cite{TerminalPressure}. The terminal measure and pressure
estimates of the latter then enter the quantitative response theorem
of \cite{ReferenceResponse}. In \cref{sec:ordered-constants} we
first verify the construction and terminal scalar conditions.
We then apply the response theorem and the volume-transfer estimates
to choose the probe-dependent interval.
